\subsection{Coverage of the compact unitary Hecke space}
\label{sec:u3-coverage}

The following lemma separates the representation-theoretic coverage from
both the finite Hecke calculation and the source's Galois-to-automorphic
normalization. Its proof uses only the proved generic, pure-inner-form case
of the global multiplicity formula.

\begin{theorem}[K17-U3-COVERAGE]\label{thm:u3-coverage}
Let $E=\mathbf Q(\sqrt{-3})$, let $G$ be the unitary group of the standard
positive definite Hermitian form on $E^3$, and put
\[
 K=\prod_q K_q,\qquad
 K_q=\operatorname{Stab}_{G(\mathbf Q_q)}
       \bigl((O_E\otimes\mathbf Z_q)^3\bigr).
\]
Suppose that $\Pi$ is a conjugate-self-dual cuspidal automorphic
representation of $\mathrm{GL}_3(\mathbf A_E)$ in unitary normalization,
unramified at every finite place, whose archimedean Langlands parameter is
\[
 z\longmapsto
 \operatorname{diag}\bigl((z/\bar z)^6,1,(z/\bar z)^{-6}\bigr).
\]
Then a discrete automorphic representation $\pi$ of $G(\mathbf A_\mathbf Q)$
exists with stable base change $\Pi$, with $\pi_f^K\ne0$, and with
$\pi_\infty$ the compact irreducible representation of highest weight
$(5,0,-5)$. Its automorphic multiplicity is one. In particular, the Hecke
system of every such $\Pi$ occurs in the algebraic modular-form space for
this compact group, this full lattice level, and this coefficient
representation.
\end{theorem}

\begin{proof}
We first specify the global theorem and its hypotheses. KMSW
\cite[Theorem~5.0.5 and its preceding paragraph, p.~206]{K17KMSW2014}
prove the full discrete multiplicity formula for generic parameters and
pure inner twists. Their \S0.3.3, p.~18, identifies the pure inner forms
as the unitary groups of Hermitian spaces over $E$, rather than over a
nontrivial division algebra. Thus $G$ satisfies this hypothesis. It is
quasi-split at every finite place, since an odd-dimensional unitary group
over a nonarchimedean local field has a unique isomorphism class, and is
compact at infinity.

By \cite[Propositions~1.3.1 and~1.3.3, pp.~67, 70--71]{K17KMSW2014},
$\Pi$ defines a simple generic global parameter with a uniquely determined
sign, and that parameter localizes at every place. The local parity
criterion in \cite[\S1.2.3 and Lemma~1.2.5, pp.~59--61]{K17KMSW2014}
forces the sign at infinity to be positive: each of the three distinct
characters $(z/\bar z)^6,1,(z/\bar z)^{-6}$ is conjugate-orthogonal,
and its multiplicity is one. Indeed an invariant conjugate pairing cannot
pair different such characters, while each character has even exponent
difference. Consequently the global parameter belongs to the standard
untwisted base-change embedding for $U(3)$.

The reduced component group of a simple generic parameter is trivial
\cite[\S1.3.4--1.3.5, pp.~69--70]{K17KMSW2014}. Hence the multiplicity
condition imposes no parity constraint on a choice of local packet members;
the unreduced central contribution disappears by the product formula in
\cite[\S1.7, pp.~88--89]{K17KMSW2014}. At infinity the displayed parameter
is discrete and relevant to the compact inner form
\cite[Remark~A.11.8, p.~236]{K17BCFamilies2009}. Its compact highest weight
is $(5,0,-5)$, because adding $\rho=(1,0,-1)$ gives the Harish--Chandra
parameter $(6,0,-6)$.

We construct fixed vectors at the finite places. The central character
$\omega_\Pi$ is unramified at every finite place and trivial at infinity.
It therefore factors through the class group of $E$. Since $O_E$ is
Euclidean, that class group is trivial, so $\omega_\Pi=1$.

At a split prime, the claim is the ordinary spherical assertion for
$\mathrm{GL}_3$. At an inert unramified prime, $K_q$ is hyperspecial.
Inertia of $\mathbf Q_q$ equals inertia of $E_q$, so the unitary-group
parameter is unramified. The spherical member of its packet has nonzero
$K_q$-invariants, by the unramified Satake correspondence and
\cite[Proposition~1.5.2(5), p.~79]{K17KMSW2014}.
Uniqueness of the unitary-group parameter above its standard base change
is \cite[Lemma~1.2.5]{K17KMSW2014}. This argument includes the inert
prime two.

For the ramified prime three we give a principal-series argument which
keeps track of the full lattice stabilizer. Write $F=\mathbf Q_3$ and
$L=E_3$. Adler--Lansky \cite[\S4.1, pp.~10--11]{K17AdlerLansky2010}
use the unitary group with Gram matrix
\[
 \Phi=\begin{pmatrix}0&0&1\\0&-1&0\\1&0&0\end{pmatrix}
\]
and its standard lattice stabilizer at the vertex $y$. Our lattice is
integrally isometric to that model. Indeed Hensel's lemma gives
$a\in\mathbf Z_3^\times$ with $a^2=-2$, so $v=(1,1,a)$ is primitive
isotropic for the standard form. Choose $w$ with $h(v,w)=1$ and replace it
by $w-\tfrac12h(w,w)v$ to obtain an integral hyperbolic plane. Its
orthogonal complement is a unit line, say of norm $d$. If $B$ is the
resulting integral basis matrix, determinants give
$-d=N_{L/F}(\det B)$. Scaling the last vector by $(\det B)^{-1}$ makes
its norm $-1$, producing the matrix $\Phi$ after ordering the basis.

The full lattice stabilizer in this model has reductive quotient
$O_3(k)=SO_3(k)\times\{\pm1\}$
\cite[\S4.2.1, pp.~11--12]{K17AdlerLansky2010}. We will obtain invariants
under this full stabilizer, so no connected-parahoric enlargement is
implicit. Conjugate self-duality and $\omega_\Pi=1$ show that the
unramified Frobenius eigenvalues of $\Pi_L$ have the form
\[
 \{a_0,1,a_0^{-1}\}.
\]
Choose an unramified character $\chi$ of $L^\times$ with Frobenius value
$a_0$, using consistent reciprocity and normalized-induction conventions.
For the Borel Levi
\[
 M=\{\operatorname{diag}(\alpha,\beta,\bar\alpha^{-1}):
       \alpha\in L^\times,\ \beta\in L^1\},
\]
set $\lambda_1=\chi$ and $\lambda_2=1$ in
\cite[(4.4.1), p.~14]{K17AdlerLansky2010}. Thus
$\lambda(\operatorname{diag}(\alpha,\beta,\bar\alpha^{-1}))=
\chi(\alpha)$, and its stable torus lift has characters
$(\chi,1,\chi^{-c})=(\chi,1,\chi^{-1})$.

The lifted $\mathrm{GL}_3(L)$ principal series is irreducible and equals
$\Pi_L$: an unramified generic local representation is its irreducible
unramified principal series, whereas at an exceptional reducibility ratio
the spherical constituent is nongeneric. Local components of a cuspidal
$\mathrm{GL}_3$ representation are generic.
The explicit unitary reducibility list
\cite[\S4.4, pp.~14--15]{K17AdlerLansky2010}
then gives only two possibilities. The induction
$I=\operatorname{Ind}_B^{G(F)}\lambda$ is irreducible, or $\chi$ is the
nontrivial unramified quadratic character and $I$ has two constituents
forming a single $L$-packet. The case $\chi=|\cdot|_L^{\pm1}$ is
excluded by irreducibility of the lifted induction. The case
$\chi|_{F^\times}=\omega_{L/F}|\cdot|_F^{\pm1}$ is excluded because
its right side is ramified and its left side is unramified.
By \cite[Proposition~6.1(1),(4), p.~21, and
Proposition~7.1(i),(iv), p.~28]{K17AdlerLansky2010}, the single packet
just described has stable base change $\Pi_L$.
These principal-series statements assume only characteristic zero and odd
residual characteristic; the exclusions $q>3$ and $q\ne7$ in the paper's
separate supercuspidal theorem are not used.

The full lattice stabilizer $K_3$ is special maximal. Iwasawa decomposition
$G(F)=B(F)K_3$, together with triviality of $\lambda$ and the modulus on
$B\cap K_3$, gives $\dim I^{K_3}=1$. Compact-open invariants are exact
over $\mathbf C$, by averaging. Hence an irreducible constituent of $I$
has a nonzero $K_3$-fixed vector. It is in the required local packet by the
preceding calculation.

We now choose these spherical members at all finite places and the compact
member at infinity. They form a restricted tensor product in the global
packet for $\Pi$. Its reduced component group is trivial, and $G$ is a
pure inner form. The proved generic case of
\cite[Theorem~5.0.5]{K17KMSW2014} therefore puts this tensor product in
the discrete spectrum with multiplicity one. The standard description of
automorphic forms on a group compact at infinity identifies its fixed
infinity type and finite-level invariants with the asserted algebraic
modular-form space.
\end{proof}

\paragraph{Scope of the coverage lemma.}
The lemma concerns every representation satisfying its stated automorphic
hypotheses. Applying it to the recovered $e=6,N=2$ source branch requires
the separately checked Weil--Deligne, twist, and coefficient normalizations.
It neither computes the compact group's class number nor determines the
dimension or Hecke eigenvalues of the finite coefficient space. Those
remain separate exact certificates. In particular the lemma alone makes
no assertion about $K_{12}(\mathbf Z)$.
